\begin{summarybox}
	\textbf{\textcolor{CSummary}{一句话核心}}

	\textbf{定义域对称是奇偶性的前提，关系式是判定的依据。}

	\medskip
	\textbf{\textcolor{CSummary}{使用条件}}
	\begin{itemize}[leftmargin=*, itemsep=0.5em, topsep=0.4em]
		\item \textbf{条件1（定义域对称）：} 函数定义域必须关于原点对称。
		\item \textbf{条件2（关系式成立）：} 偶函数：$f(-x)=f(x)$；奇函数：$f(-x)=-f(x)$。
	\end{itemize}

	\medskip
	\textbf{\textcolor{CSummary}{关键提醒}}
	\begin{itemize}[leftmargin=*, itemsep=0.5em, topsep=0.4em]
		\item \textbf{易错点（忽略对称性）：} 直接计算$f(-x)$而不检查定义域，可能误判。
		\item \textbf{检查项（验证完整）：} 先看定义域是否对称，再算$f(-x)$并与$f(x)$比较。
	\end{itemize}
\end{summarybox}
